\chapter{INTRODUCTION}

% =========================================================
\section{Background of the Study}
% =========================================================
The rapid proliferation of digital media on the internet has led to an increasing demand for robust copyright protection and secure data communication. Steganography, the art and science of hiding information within other media in such a way that no one, apart from the intended recipient, usually expects existing of the message, has emerged as a key solution. Unlike cryptography, which scrambles data to make it unreadable, steganography conceals the very existence of the data. 

In the realm of digital images, the RGB color space offers a vast landscape for data embedding. Recent advancements, such as those by Liu and Su \cite{liu2025color}, have explored combining Pixel Value Differencing (PVD) with edge detection to enhance embedding capacity while maintaining visual imperceptibility. However, traditional approaches often rely on sequential embedding, where data is hidden in one channel (e.g., Red) before moving to the next (e.g., Green, then Blue). This sequential dependency can create bottlenecks in capacity and introduce complex inter-channel artifacts that are susceptible to steganalysis.

This research introduces a "Multiple Embedding Steganography" (MES) framework that shifts the paradigm from sequential to parallel embedding. By treating the R, G, and B channels as independent carriers processed simultaneously, we aim to significantly unlock higher payload capacities required for embedding large data, such as full-size secret images, without compromising the visual integrity of the cover image.

% =========================================================
\section{Statement of the Problem}
% =========================================================
Existing high-capacity steganographic methods, particularly those operating in the RGB domain, face the following challenges:
\begin{enumerate}
    \item \textbf{Sequential Dependency Limitation:} Most current algorithms embed data sequentially across color channels. This creates a dependency where the embedding capacity and modification in one channel may constrain or affect the others, limiting the overall payload.
    \item \textbf{Capacity vs. Imperceptibility Trade-off:} Achieving high embedding capacity often results in significant distortion, lowering the Peak Signal-to-Noise Ratio (PSNR) and making the stego-image visually detectable.
    \item \textbf{Limited Payload for Large Data:} Current methods struggle to embed large payloads, such as a full-resolution image within another image of the same size, without causing noticeable degradation.
\end{enumerate}

% =========================================================
\section{Research Objectives}
% =========================================================
The primary objective of this study is to develop a Multiple Embedding Steganography (MES) framework for RGB images that utilizes a parallel embedding architecture.
Specific objectives include:
\begin{enumerate}
    \item To design a parallel embedding algorithm that processes Red, Green, and Blue channels independently to maximize spatial utilization.
    \item To implement an adaptive embedding strength mechanism based on edge detection to preserve visual quality in smooth regions.
    \item To achieve an embedding capacity greater than 3.0 bits per pixel (bpp).
    \item To maintain a high visual quality with a PSNR value above 30 dB.
\end{enumerate}

% =========================================================
\section{Scope and Delimitation}
% =========================================================
\textbf{Scope:}
\begin{itemize}
    \item The study focuses on digital images in the RGB color space.
    \item The proposed method utilizes Pixel Value Differencing (PVD) and Edge Detection techniques.
    \item Performance evaluation will be based on Embedding Capacity (bpp), PSNR, and SSIM metrics.
\end{itemize}

\textbf{Delimitations:}
\begin{itemize}
    \item This research does not cover robustness against active attacks (e.g., crop, rotation, compression).
    \item The study is limited to spatial domain embedding and does not explore frequency domain techniques (e.g., DCT/DWT).
\end{itemize}

% =========================================================
\section{Significance of the Study}
% =========================================================
This research contributes to the field of information security by demonstrating that parallel processing in the color domain can significantly enhance steganographic capacity. The findings will be beneficial for applications requiring the covert transmission of large data files, such as medical imaging privacy and secure military communication, providing a method that offers both high capacity and high imperceptibility.
